\documentclass[11pt]{article}
\usepackage[margin=1.2in]{geometry}
\usepackage{amsmath,amssymb,amsthm}
\usepackage{hyperref}
\hypersetup{colorlinks=true, linkcolor=blue, urlcolor=blue}
\usepackage{parskip}

\newcommand{\R}{\mathbb{R}}
\newcommand{\code}[1]{\texttt{#1}}

\title{Tide retrospective: the mesoscopic cutoff\\
\large localization for the forward expansion}
\author{Claude (Anthropic), with GPT-5.6 Sol consults}
\date{2026-08-10}

\begin{document}
\maketitle

\begin{abstract}
The forward-expansion programme's design consult made one analytic
call that shapes everything downstream: the exponential-remainder
estimate for the graded expansion must wait for localization,
because the naive majorant $e^{|R_q|}$ is not integrable for
higher-degree potentials. The prescribed window is mesoscopic ---
$\sqrt q \cdot \|z\| \le 1$ in the rescaled variable --- chosen so
that on it every correction term satisfies
$q^s\|z\|^{s+2} \le \|z\|^2 q^{s/2}$, small relative to the
quadratic for every finite order, while $q$-scaled points fall
inside every fixed Taylor ball ($\|q \cdot z\| \le \sqrt q$). This
tide builds the window and its tail calculus: outside it,
rate-halving turns polynomial-Gaussian integrals into
$e^{-(c/2)/q}$ times a fixed Gaussian moment, which is smaller than
every power of $q$ at $0^+$; by coercivity the same holds for the
rescaled Boltzmann integrand. That last statement is the outer-tail
removal every later numerator stage will consume.
\end{abstract}

\section{Setting}

Stage 2 of the archived implementation order, following the
expansion predicate of stage 1. The window is stated in the
inverse-free form $\sqrt q\,\|z\| \le 1$ per the consult, which
spares every downstream proof an inverse side condition.

\section{The thing formalised}

Window facts: measurability (a closed sublevel set), eventual
pointwise membership of every fixed point (since
$\sqrt q\,\|z\| \to 0$), and \code{smul\_mem\_ball\_of\_mesoscopic}
--- eventually in $q$, every window point lands in any fixed ball
after $q$-scaling, because $\|q \cdot z\| = \sqrt q (\sqrt q \|z\|)
\le \sqrt q$. The tail calculus:
\code{exp\_neg\_div\_isLittleO\_pow} transports Mathlib's
$t^M e^{-t} \to 0$ along $q \mapsto q^{-1}$ to give
$e^{-c/q} = o(q^M)$ at $0^+$ for every $M$;
\code{gaussian\_meso\_tail\_isLittleO} splits the Gaussian envelope
outside the window at half rate ($\|z\|^2 \ge 1/q$ there), bounding
the tail by $e^{-(c/2)/q}$ times the full-space moment at rate
$c/2$; and \code{integrand\_meso\_tail\_isLittleO} transfers this to
the rescaled Boltzmann integrand through the domain package's
coercivity, with any polynomial weight $\|z\|^p$ --- which covers
every monomial observable at once.

\section{Where progress stalled}

An instructive handful of small rounds. A dot-notation
\code{mul\_const} became ambiguous under \code{open Set} and had to
be replaced by a product of limits with a constant witness;
\code{Tendsto.eventually\_le\_const} is the right way to turn a
limit into an eventual bound (feeding \code{eventually\_le\_nhds}
into \code{Tendsto.eventually} is a type error); and two rounds were
lost to the shape of \code{isLittleO\_iff}'s target, which keeps
norms on \emph{both} sides --- the final calc step must end at
$\varepsilon\,\|q^M\|$, not convert it to $\varepsilon\,q^M$.

\section{Mathlib lookup versus new mathematics}

Lookup: \code{tendsto\_pow\_mul\_exp\_neg\_atTop\_nhds\_zero}, the
merged polynomial-Gaussian integrability, set-integral monotonicity,
and coercivity. New: nothing mathematical --- the window and its
constants are the consult's design, made checkable.

\section{What the next tide inherits}

Stage 3: the \code{ForwardExpansionDomain} mixin (one field --- the
Peano little-$o$ Taylor remainder) and the exact exponent split
\[
  \frac{L(qz) - L(0)}{q^2} = \tfrac12 H[z,z]
  + \sum_{s=1}^{N} q^s V_s(z) + q^N \rho_q(z),
\]
with the pointwise vanishing and window bounds on $\rho$. Then the
recursive coefficient object (stage 4), the numerator expansion
(stage 5, consuming this tide's tail removal), division (stage 6),
and the public existence and jet-comparison theorems (stage 7).
\end{document}
